\documentclass[hidelinks]{article}
\usepackage[a4paper, total={7in, 10in}]{geometry}
\usepackage[dvipsnames]{xcolor}
\usepackage{amsmath, amssymb, amsthm}
\usepackage{tikz}
\usepackage{tkz-euclide}
\usepackage[unicode]{hyperref}
\usepackage[all]{hypcap}
\usepackage{fancyhdr}
\usepackage{listings}
\usepackage{subcaption}
\usepackage{enumitem}
\usepackage{graphicx}

\lstset{basicstyle=\ttfamily\small,columns=fullflexible,frame=single,breaklines=true,showstringspaces=false}
\usetikzlibrary{angles,calc,decorations.pathreplacing}

\title{\textbf{Laplace's Equation Test}}
\author{Diego Juarez}
\date{March 25th, 2026}

\pagestyle{fancy}
\fancyhf{}
\fancyhead[L]{Laplace's Equation Test}
\fancyhead[R]{Jackson-style Practice}
\fancyfoot[C]{\thepage}

\begin{document}
\hypersetup{bookmarksnumbered=true}
\pagecolor{black}
\color{white}
\maketitle

\begin{Large}
\tableofcontents
\end{Large}
\pagebreak

\section{Source Notebook}
This problem set follows the notebook
\url{../notebooks/Laplace's equation.ipynb}
and emphasizes harmonic functions, separation of variables, and uniqueness.

\section{Foundations}
\subsection{Problem 1: Laplace Equation from Poisson Equation}
Starting from Poisson's equation, explain why in a charge-free region the electrostatic potential satisfies
\[
\nabla^2\Phi=0.
\]
State the physical meaning of a source-free region.

\subsection{Problem 2: Mean Value Property}
State the mean value property for harmonic functions. Then explain why it immediately implies that a harmonic function cannot have an isolated interior maximum or minimum unless it is constant.

\subsection{Problem 3: Uniqueness}
Prove the Dirichlet uniqueness theorem for Laplace's equation in electrostatics. Keep the proof at the level of an energy integral argument.

\section{One- and Two-Dimensional Problems}
\subsection{Problem 4: One-Dimensional Laplace Equation}
Solve
\[
\frac{d^2\Phi}{dx^2}=0
\]
on the interval $0<x<L$ with boundary conditions $\Phi(0)=V_0$ and $\Phi(L)=V_1$. Interpret the electric field physically.

\subsection{Problem 5: Parallel Plates}
Two infinite conducting plates are located at $x=0$ and $x=d$, held at potentials $0$ and $V_0$ respectively.
\begin{enumerate}[label=\alph*)]
\item Find the potential between the plates.
\item Find the electric field.
\item Explain why fringing fields are absent in this idealization.
\end{enumerate}

\subsection{Problem 6: Rectangular Region}
Solve Laplace's equation in the rectangle $0<x<a$, $0<y<b$ with boundary conditions
\[
\Phi(0,y)=0,\quad \Phi(a,y)=0,\quad \Phi(x,0)=0,\quad \Phi(x,b)=V_0\sin\left(\frac{\pi x}{a}\right).
\]
Use separation of variables and determine the full solution.

\subsection{Problem 7: General Fourier-Series Boundary Data}
Repeat the previous problem when the boundary value at $y=b$ is an arbitrary function $f(x)$ expanded in a sine series. Write the solution as a series without evaluating the coefficients explicitly.

\section{Symmetry-Based Solutions}
\subsection{Problem 8: Spherical Symmetry}
Assume a radially symmetric solution $\Phi(r)$ in a charge-free region.
\begin{enumerate}[label=\alph*)]
\item Reduce Laplace's equation to an ordinary differential equation.
\item Solve for $\Phi(r)$.
\item Explain the physical meaning of each integration constant.
\end{enumerate}

\subsection{Problem 9: Exterior of a Conducting Sphere}
A conducting sphere of radius $R$ is held at potential $V_0$.
\begin{enumerate}[label=\alph*)]
\item Solve for the exterior potential assuming $\Phi(\infty)=0$.
\item Find the electric field outside the sphere.
\item Compute the surface charge density from the discontinuity of the normal electric field.
\end{enumerate}

\subsection{Problem 10: Cylindrical Symmetry}
Assume $\Phi=\Phi(s)$ depends only on cylindrical radius $s$.
\begin{enumerate}[label=\alph*)]
\item Derive the radial form of Laplace's equation in cylindrical coordinates.
\item Solve it.
\item Explain why logarithms appear in two-dimensional electrostatics.
\end{enumerate}

\section{Properties of Harmonic Functions}
\subsection{Problem 11: No Interior Extrema}
Suppose $\Phi$ is harmonic in a bounded domain. Give a proof, based on either the mean value property or the maximum principle, that $\Phi$ cannot attain a strict interior maximum unless it is constant.

\subsection{Problem 12: Consequence for Conductors}
Explain why a conductor in electrostatic equilibrium must have constant potential throughout its interior. Then connect this result to Laplace's equation in the empty region outside the conductor.

\section{Boundary Value Perspective}
\subsection{Problem 13: Mixed Boundary Conditions}
Discuss the difference between Dirichlet and Neumann boundary conditions for Laplace's equation. Give one electrostatic example of each.

\subsection{Problem 14: Changing Boundary Data}
Explain why changing the potential on only one boundary segment can alter the solution everywhere in a connected region. Relate your answer to harmonicity and nonlocality.

\section{Challenge Problems}
\subsection{Problem 15: Harmonic Polynomial Test}
Determine which of the following are harmonic:
\[
x^2-y^2,\qquad xy,\qquad x^2+y^2,\qquad z^3-3zx^2.
\]
Compute the Laplacian explicitly in each case.

\subsection{Problem 16: Annular Region}
Solve Laplace's equation in the annulus $a<r<b$ for a radially symmetric potential with boundary values $\Phi(a)=V_a$ and $\Phi(b)=V_b$. Compare the result with the corresponding three-dimensional spherically symmetric problem.

\end{document}
